\section*{Chapitre \ref{ch:rc:mortgage-modelling} --- Modélisation des crédits immobiliers}

\begin{solution}{exo:rc:mortgage-modelling:1}
Le taux du prêt immobilier vaut $3.80+1.75 = 5.55~\%$~; l'incitation $6.00-5.55 = 0.45$ point de pourcentage,
juste en dessous du centre de la courbe en S.
\end{solution}

\begin{solution}{exo:rc:mortgage-modelling:2}
$181-146 = 34$ points de base (35 avant arrondi des deux spreads)~: la part du rendement du pool
au-dessus de la courbe qui rémunère les options de remboursement anticipé des emprunteurs lorsque les taux sont
volatils, mesurée sous forme de spread.
\end{solution}

\begin{solution}{exo:rc:mortgage-modelling:3}
Des taux plus élevés ralentissent le refinancement, si bien que le capital restant dû subsiste plus longtemps~;
les intérêts supplémentaires l'emportent sur l'effet d'une actualisation plus forte.
\end{solution}

\begin{solution}{exo:rc:mortgage-modelling:4}
Le remboursement anticipé dépend du chemin suivi (l'\omterm{def:rc:mortgage-modelling:scurve}{essoufflement} (burnout), le capital restant dû), et non seulement du taux
courant~: deux chemins atteignant le même nœud portent des pools différents. Un arbre recombinant oublie
le chemin~; la simulation le conserve.
\end{solution}

\begin{solution}{exo:rc:mortgage-modelling:5}
Une baisse de 100 points de base rapporte environ 4.54, une hausse fait perdre environ 5.50~: les pertes dépassent les
gains, c'est l'asymétrie de la convexité négative.
\end{solution}

\begin{solution}{exo:rc:mortgage-modelling:6}
Sans \omterm{def:rc:mortgage-modelling:scurve}{essoufflement}, les emprunteurs continuent de se refinancer au taux que donne la courbe en S, si bien que le pool
à prime (coupon supérieur à celui du marché) rembourse plus vite son principal au pair, ce qui réduit plus tôt
la durée pendant laquelle il verse le coupon au-dessus du marché~: sa valeur baisse, à 99.20.
\end{solution}

\begin{solution}{exo:rc:mortgage-modelling:7}
DV01 avant $4.99\times100/100\times10\text{ milliards}\times10^{-4} = \text{USD}~4.99$ million par
point de base, après $6.17\times94.50/100\times\ldots = 5.83$ million~; le surcroît de USD~844\,025 par
point de base divisé par le swap à dix ans (USD~796\,282 par milliard de notionnel), donne USD~1.06 milliard de
swaps payeurs.
\end{solution}

\begin{solution}{exo:rc:mortgage-modelling:8}
La couverture en duration ne supprime que le risque du premier ordre~: le portefeuille est vendeur de convexité (et vendeur de
volatilité au travers des options des emprunteurs), si bien qu'il perd sur les fortes variations, dans un sens comme dans l'autre
entre deux rééquilibrages, paie la fourchette acheteur--vendeur à chaque opération, et porte le risque du \omterm{def:rc:mortgage-modelling:model}{modèle de remboursement anticipé} et
du spread hypothécaire, que nulle couverture de taux ne couvre.
\end{solution}

\begin{solution}{pb:rc:mortgage-modelling:1}
\textbf{1.} 146, 181 et 34 points de base.
\textbf{2.} 4.99 ans et $-123$.
\textbf{3.} USD~4.99 million par point de base.
\textbf{4.} Intérêts 25.26, principal 74.74 pour 100.
\textbf{5.} Le pool~: il est vendeur des options de remboursement anticipé des emprunteurs, vendues implicitement par
le prêt qui leur est consenti à taux fixe, avec un droit de remboursement anticipé libre.
\textbf{6.} 94.50.
\textbf{7.} 6.17 ans~; USD~5.83 million par point de base.
\textbf{8.} Des taux plus élevés sortent le pool de la \omterm{def:rc:mortgage-modelling:scurve}{courbe en S du remboursement anticipé}~: les remboursements anticipés ralentissent et
le principal revient plus tard.
\textbf{9.} USD~1.06 milliard de swaps payeurs à dix ans.
\textbf{10.} Leurs swaps payeurs poussent les taux de swap encore à la hausse, ce qui rallonge de nouveau les pools,
d'où un besoin accru de couverture~: une rétroaction que l'étude de la Federal Reserve a relevée dans les données.
\textbf{11.} IO 24.00 à $-25$ et 26.61 à $+25$~; PO 77.21 et 72.10.
\textbf{12.} L'IO, dont la valeur croît avec les taux, compense en partie la convexité négative du PO
(et du pool).
\textbf{13.} Les dépôts dont le taux se réajuste lentement valent davantage quand les taux montent~; un IO se comporte de la
même façon et peut couvrir le bilan de la banque (chapitre~24).
\textbf{14.} Valeur~: le principal arrive plus tard que prévu, pour le même montant au pair.
\textbf{15.} L'IO~: sa valeur dépend presque entièrement de la durée pendant laquelle le capital restant dû subsiste.
\textbf{16.} Les swaptions permettraient de racheter de la convexité et coûteraient la prime de l'option (le \omterm{def:rc:mortgage-modelling:optioncost}{coût
de l'option} du pool sert d'indice de référence)~; beaucoup de détenteurs combinent les deux.
\textbf{17.} La duration s'allonge progressivement à mesure que les taux montent, et les flux de rééquilibrage s'étalent sur des semaines~;
la volatilité augmente et alimente à son tour le mouvement, et l'étude de la Federal Reserve constate que les effets persistent pendant des mois.
\textbf{18.} Un détenteur qui ne se couvre pas supprime une partie de la rétroaction~; celui qui cesse
d'acheter (laisse ses positions s'amortir) reporte la couverture sur les détenteurs privés, et peut l'alourdir.
\textbf{19.} Résultat nommé~: \emph{l'événement de convexité}~: sur USD~10 milliards du pool, une hausse de 100 points de base
fait passer le DV01 de USD~4.99 million à 5.83 million par point de base, et le rééquilibrage mobilise
USD~1.06 milliard de swaps payeurs à dix ans.
\textbf{20.} Parce que les emprunteurs cessent de se refinancer quand les taux montent, précisément au moment où le détenteur
voudrait récupérer ses fonds.
\end{solution}

\begin{solution}{iq:rc:mortgage-modelling:1}
Les emprunteurs peuvent rembourser au pair à tout moment~: le gain de prix du titre est plafonné lorsque
les taux baissent (il se raccourcit), et ses pertes croissent quand les taux montent (il s'allonge). C'est une obligation
diminuée d'une option d'achat détenue par les emprunteurs.
\iqlookfor{l'option de remboursement anticipé incorporée et les variations de prix asymétriques.}
\end{solution}

\begin{solution}{iq:rc:mortgage-modelling:2}
\omterm{def:rc:mortgage-modelling:turnover}{Rotation du parc immobilier} (ventes de logements, défauts, remboursements partiels~; pilotée par l'âge, peu sensible aux taux)~;
refinancement (une courbe en S de l'\omterm{def:rc:mortgage-modelling:turnover}{incitation au refinancement}, WAC moins taux du prêt, atténuée par l'\omterm{def:rc:mortgage-modelling:scurve}{essoufflement}
et les caractéristiques du pool)~; le taux du prêt découle des taux de marché par le
\omterm{def:rc:mortgage-modelling:cc}{spread primaire--secondaire}. Le refinancement et le taux du prêt dépendent des taux.
\iqlookfor{les composantes et celles qui dépendent des taux.}
\end{solution}

\begin{solution}{iq:rc:mortgage-modelling:3}
Simuler des trajectoires de taux à partir d'un modèle de structure par terme calibré~; faire tourner le \omterm{def:rc:mortgage-modelling:model}{modèle de remboursement anticipé} et
les flux de trésorerie du pool sur chaque trajectoire~; chercher le spread constant au-dessus du taux court
de la trajectoire qui rend le flux actualisé moyen égal au prix. Monte Carlo, parce que
le remboursement anticipé dépend du chemin suivi.
\iqlookfor{trajectoires, flux de trésorerie, recherche de racine pour le spread, et dépendance au chemin.}
\end{solution}

\begin{solution}{iq:rc:mortgage-modelling:4}
Leurs pools s'allongent~; pour conserver la duration, ils paient le taux fixe dans des swaps ou vendent des bons du Trésor américain, ce qui alimente
la vente~; la volatilité implicite augmente à mesure que les hedgers achètent des options~; le mouvement peut aller au-delà
avant que les flux ne s'épuisent.
\iqlookfor{flux de couverture dans le sens du mouvement.}
\end{solution}

\begin{solution}{iq:rc:mortgage-modelling:5}
Faire tourner plusieurs \omterm{def:rc:mortgage-modelling:model}{modèles de remboursement anticipé} ou des paramètres perturbés (multiplicateurs de vitesse, décalages de la courbe en S,
\omterm{def:rc:mortgage-modelling:scurve}{essoufflement}), calculer le spread ajusté des options (OAS) et le risque dans chaque cas, présenter les durations du
remboursement anticipé, et provisionner la dispersion des valeurs~; tester a posteriori le modèle contre les vitesses
réalisées par cohorte.
\iqlookfor{dispersion des scénarios et back-testing.}
\end{solution}

\begin{solution}{iq:rc:mortgage-modelling:6}
Nombres aléatoires non communs entre les calculs perturbés (une nouvelle graine à chaque fois), trop peu de trajectoires,
ensembles de trajectoires de taille variable, courbe ou volatilité qui bougent, et la
perturbation aux différences finies qui interagit avec une courbe en S abrupte~; fixer les graines et les trajectoires, utiliser les variables antithétiques, et vérifier la
taille de la perturbation.
\iqlookfor{nombres aléatoires communs et diagnostic du bruit.}
\end{solution}
